% !PRINTED 268
\weberpage{268}
In der Gruppe $\Lambda_p$ ist die Normalform $S$ für die Substitutionen
 dritter Art enthalten.

Nehmen wir nun eine Substitution $A$ von der dritten Art in
der Gruppe $\Lambda_p$ an, so können wir sie, wie jetzt noch nachgewiesen
 werden soll, durch Transformation mit einer Substitution
 $U$, die selbst der Gruppe $\Lambda_p$ angehört, in die Normalform
 transformiren. Denn setzen wir
\weq{8}{A = \begin{pmatrix} \alpha, & \beta \\ -N\beta', & \alpha' \end{pmatrix}, \quad U = \begin{pmatrix} \lambda, & \mu \\ -N\mu', & \lambda' \end{pmatrix},}
\weq{9}{\alpha\alpha'+N\beta\beta' = 1, \qquad \lambda\lambda'+N\mu\mu' = 1,}
worin $\lambda,\lambda'$ und $\mu,\mu'$ conjugirte Paare sind, so erhalten die
Gleichungen (3), (4), §.\,68 die Form:
\weq{10}{(\sigma-\alpha)(\sigma-\alpha')+N\beta\beta' = \sigma^2-\sigma(\alpha+\alpha')+1 = 0,}
\weq{11}{N\,\dfrac{\mu'}{\lambda} = \dfrac{\alpha-\sigma}{\beta}, \qquad \dfrac{\lambda'}{\mu} = -\dfrac{\alpha-\sigma^{-1}}{\beta},}
und die Gleichung (10) muss, da $A$ von der dritten Art sein
soll, zwei zu einander reciproke, conjugirt imaginäre Wurzeln
haben.

Aus der zweiten Gleichung (11) folgt aber durch Uebergang
zu den conjugirt imaginären Zahlen
\[
\dfrac{\lambda}{\mu'} = -\dfrac{\alpha'-\sigma}{\beta'},
\]
und dies ist nach (10) eine Folge der ersten Gleichung (11).
Setzen wir also
\[
N\mu' = \varkappa\,(\alpha-\sigma), \qquad N\mu = \varkappa'\,(\alpha'-\sigma^{-1})
\]
\[
\lambda = \varkappa\beta, \qquad \lambda' = \varkappa'\beta',
\]
worin $\varkappa,\varkappa'$ zwei conjugirt imaginäre Zahlen sind, so sind die
Gleichungen (11) befriedigt, und man kann dann $\varkappa\varkappa'$ noch so
bestimmen, dass $\lambda\lambda'+N\mu\mu' = 1$ wird.

Daraus ergiebt sich noch ein wichtiges Resultat. Bezeichnen
wir mit $\gamma$ eine primitive Wurzel des Körpers $\mathfrak{E}_{2,p}$, so können
wir für die Substitutionen der Normalform für die zweite und
dritte Art die Ausdrücke annehmen (§.\,68):
\[
S_2^r = \begin{pmatrix} \gamma^{(p+1)r}, & 0 \\ 0, & \gamma^{-(p+1)r} \end{pmatrix}, \qquad S_3^r = \begin{pmatrix} \gamma^{(p-1)r}, & 0 \\ 0, & \gamma^{-(p-1)r} \end{pmatrix}.
\]

Ist nun $A$ eine Substitution zweiter Art aus $L_p$ und $A$ eine
%CONT>>
